\documentclass[11pt]{article}
\usepackage[margin=1in]{geometry}
\usepackage{amsmath, amssymb, amsthm}
\usepackage{hyperref}
\usepackage{enumitem}

\title{Math - Test Sections Questions: Free Mechanical Vibrations}
\author{Curated Problem Sets}
\date{}

\begin{document}

\maketitle
\tableofcontents
\newpage

% ---------------------------------------------------------
\section{Concept Checklist \& Question Mapping}
% ---------------------------------------------------------
This section outlines all the core concepts, problem types, and tricks found in the provided text on Free Mechanical Vibrations. Each concept is mapped to the corresponding problems in the Problem Sets to ensure comprehensive coverage.

\begin{itemize}
    \item \textbf{Concept 1: Undamped Free Vibrations (IVP \& Equation of Motion)} - Solving $my'' + ky = 0$. (Problems: 1, 4, 28)
    \item \textbf{Concept 2: Amplitude, Period, Frequency, and Phase Angle} - Converting to $A\sin(\omega t + \phi)$ and finding $P = 2\pi/\omega$, $f = \omega/2\pi$. Checking quadrants for $\phi$. (Problems: 1, 3, 5, 28)
    \item \textbf{Concept 3: Time to First Equilibrium (Undamped)} - Setting $y(t) = 0$ and solving for the first $t > 0$. (Problems: 2)
    \item \textbf{Concept 4: Damped Free Vibrations - Overdamped} - Solving $my'' + by' + ky = 0$ where $b^2 > 4mk$. (Problems: 6, 15, 19)
    \item \textbf{Concept 5: Damped Free Vibrations - Critically Damped} - Solving where $b^2 = 4mk$, using repeated root template $(C_1 + C_2t)e^{rt}$. (Problems: 8, 14, 18)
    \item \textbf{Concept 6: Damped Free Vibrations - Underdamped} - Solving where $b^2 < 4mk$, using complex roots template. (Problems: 10, 16, 29)
    \item \textbf{Concept 7: Damping Factor, Quasiperiod, Quasifrequency} - Identifying $Ae^{\alpha t}$, $2\pi/\beta$, and $\beta/2\pi$. (Problems: 11, 12)
    \item \textbf{Concept 8: Maximum Displacement (Optimization)} - Setting $y'(t) = 0$ to find max displacement in damped/undamped systems. (Problems: 9, 17, 19)
    \item \textbf{Concept 9: Time to First Equilibrium (Damped)} - Setting $y(t) = 0$ for damped systems and solving for $t$. (Problems: 7, 13, 18)
    \item \textbf{Concept 10: Theoretical Proofs \& Limits} - Algebraic proofs and limits (e.g., $t \to \infty$). (Problems: 20, 30)
    \item \textbf{Concept 11: Forced Vibrations \& Resonance} - Nonhomogeneous ODEs $my'' + by' + ky = F_0\cos(\gamma t)$ and the Method of Undetermined Coefficients. (Problems: 21, 22)
    \item \textbf{Concept 12: Diagnostic Testing ("Find the Flaw")} - Identifying common pitfalls like discriminant errors, phase angle quadrant errors, missing $t$ in repeated roots, and derivative product rule errors. (Problems: 23, 24, 25, 26, 27)
\end{itemize}

\newpage

% ---------------------------------------------------------
\section{Problem Sets}
% ---------------------------------------------------------

\begin{enumerate}[label=\textbf{\arabic*.}]
    % Undamped
    \item A 1-kg mass is attached to a spring with stiffness $k = 16$ N/m. The mass is displaced 0.5 m to the left of the equilibrium point and given a velocity of 2 m/sec to the left. Neglecting damping, find the equation of motion of the mass. Convert your answer to the phase-amplitude form $y(t) = A\sin(\omega t + \phi)$.
    \item Using the equation of motion derived in Problem 1, determine how long after release the mass first passes through the equilibrium position.
    \item A 4-kg mass is attached to a spring with stiffness $k = 25$ N/m. There is no damping. Find the period and the natural frequency of the resulting motion.
    \item A 2-kg mass is attached to a spring with stiffness $k = 18$ N/m. The mass is pulled 1 m to the right of equilibrium and given an initial rightward velocity of 3 m/sec. Find the equation of motion $y(t)$.
    \item Convert the equation of motion found in Problem 4 into the phase-amplitude form $y(t) = A\sin(\omega t + \phi)$.
    
    % Overdamped
    \item A 1-kg mass is attached to a spring with stiffness $k = 4$ N/m. The damping constant is $b = 5$ N-sec/m. The mass is pulled 2 m to the right and given an initial leftward velocity of 1 m/sec. Find the equation of motion.
    \item For the system described in Problem 6, determine the time $t > 0$ when the mass first passes through the equilibrium position. If it never passes through equilibrium, prove it mathematically.
    
    % Critically Damped
    \item A 1-kg mass is attached to a spring with stiffness $k = 4$ N/m. The damping constant is $b = 4$ N-sec/m. The mass is pulled 1 m to the right and released from rest. Find the equation of motion.
    \item Find the maximum displacement from equilibrium for the mass in Problem 8.
    
    % Underdamped
    \item A 1-kg mass is attached to a spring with stiffness $k = 5$ N/m. The damping constant is $b = 2$ N-sec/m. The mass is pulled 2 m to the right and given an initial leftward velocity of 2 m/sec. Find the equation of motion.
    \item For the system in Problem 10, determine the quasiperiod and the quasifrequency of the motion.
    \item For the system in Problem 10, determine the damping factor of the motion.
    \item For the system in Problem 10, find the first time $t > 0$ that the mass passes through the equilibrium position.
    
    % Determine Damping Type
    \item A mass-spring system is governed by $3y'' + 12y' + 12y = 0$. Determine if the system is underdamped, critically damped, or overdamped.
    \item A mass-spring system is governed by $2y'' + 5y' + 3y = 0$. Determine if the system is underdamped, critically damped, or overdamped.
    \item A mass-spring system is governed by $y'' + 2y' + 10y = 0$. Determine if the system is underdamped, critically damped, or overdamped.
    
    % Max Displacement
    \item A 1-kg mass is attached to a spring with stiffness $k = 2$ N/m and a damping constant $b = 2$ N-sec/m. It is pushed from equilibrium ($y(0)=0$) with a rightward velocity of 1 m/sec. Find the maximum displacement it attains to the right.
    \item A critically damped system has $m=1$, $b=6$, and $k=9$. The mass is pulled 2 m to the right and given a severe leftward velocity of 10 m/sec. Find the exact time $t > 0$ when the mass passes through equilibrium.
    \item An overdamped system has the equation of motion $y(t) = e^{-t} - e^{-3t}$. Find the maximum displacement of the mass from equilibrium.
    
    % Theoretical
    \item Prove that for any overdamped system with equation of motion $y(t) = C_1 e^{r_1 t} + C_2 e^{r_2 t}$ (where $r_1 \neq r_2$), the mass can pass through the equilibrium position at most once.
    
    % Forced Vibrations
    \item An undamped system is subjected to an external force, governed by $y'' + 4y = 3\cos(t)$. Find the general solution for the motion of the mass.
    \item An undamped system is governed by $y'' + 16y = 5\sin(4t)$. Without solving for the exact coefficients, write down the correct form of the particular solution $y_p(t)$ and state the name of the physical phenomenon that will occur.
    
    % Find the Flaw
    \item \textbf{Find the Flaw:} A student is asked to determine the damping type for $m=2, b=4, k=4$. 
    \textit{Flawed Solution:} 1. Calculate the discriminant $b^2 + 4mk$. 2. $16 + 32 = 48$. 3. Since $48 > 0$, the system is overdamped. 
    Identify the error, explain why it is wrong in one sentence, and provide the correct final answer.
    
    \item \textbf{Find the Flaw:} A student converts $y(t) = -2\cos(t) - 2\sin(t)$ into $A\sin(t + \phi)$. 
    \textit{Flawed Solution:} 1. $A = \sqrt{(-2)^2 + (-2)^2} = \sqrt{8} = 2\sqrt{2}$. 2. $\tan \phi = \frac{-2}{-2} = 1$. 3. $\phi = \arctan(1) = \pi/4$. 4. $y(t) = 2\sqrt{2}\sin(t + \pi/4)$.
    Identify the error, explain why it is wrong in one sentence, and provide the correct final answer.
    
    \item \textbf{Find the Flaw:} A student solves the IVP $y'' + 10y' + 25y = 0$. 
    \textit{Flawed Solution:} 1. Characteristic equation $r^2 + 10r + 25 = 0$. 2. $(r+5)^2 = 0 \implies r = -5$. 3. General solution is $y(t) = C_1 e^{-5t} + C_2 e^{-5t}$.
    Identify the error, explain why it is wrong in one sentence, and provide the correct final answer.
    
    \item \textbf{Find the Flaw:} A student finds the time of maximum displacement for $y(t) = e^{-t}\sin(t)$. 
    \textit{Flawed Solution:} 1. Set $y(t) = 0$ to find the maximum. 2. $e^{-t}\sin(t) = 0$. 3. Since $e^{-t} \neq 0$, $\sin(t) = 0 \implies t = \pi$.
    Identify the error, explain why it is wrong in one sentence, and provide the correct final answer.
    
    \item \textbf{Find the Flaw:} A student calculates the velocity $y'(t)$ for $y(t) = 2e^{-3t}\cos(2t)$. 
    \textit{Flawed Solution:} 1. Derivative of $2e^{-3t}$ is $-6e^{-3t}$. 2. Derivative of $\cos(2t)$ is $-2\sin(2t)$. 3. Multiply them: $y'(t) = 12e^{-3t}\sin(2t)$.
    Identify the error, explain why it is wrong in one sentence, and provide the correct final answer.
    
    % Mixed / Advanced
    \item A 1-kg mass is attached to a spring with stiffness $k = 2$ N/m. The mass is pulled 1 m to the right and given a leftward velocity of 1 m/sec. Find the phase-amplitude form of the equation of motion.
    \item An underdamped system has $m=1, b=2, k=10$. The mass is released from equilibrium ($y(0)=0$) with an initial rightward velocity of 6 m/sec. Find the equation of motion.
    \item Consider a critically damped system with equation of motion $y(t) = (C_1 + C_2 t)e^{-rt}$, where $r = b/2m > 0$. Prove using calculus that regardless of the initial conditions (the values of $C_1$ and $C_2$), the mass will eventually return to the equilibrium position as $t \to \infty$ (i.e., $\lim_{t \to \infty} y(t) = 0$).
\end{enumerate}

\newpage

% ---------------------------------------------------------
\section{Detailed Solutions}
% ---------------------------------------------------------

\begin{enumerate}[label=\textbf{\arabic*.}]
    
    % Solution 1
    \item \textbf{Set up the IVP:} $m=1, k=16, b=0$. $y(0) = -0.5, y'(0) = -2$. \\
    $y'' + 16y = 0 \implies r^2 + 16 = 0 \implies r = \pm 4i$. \\
    $y(t) = C_1 \cos(4t) + C_2 \sin(4t)$. \\
    \textbf{Apply Initial Conditions:} $y(0) = C_1 = -0.5$. \\
    $y'(t) = -4C_1 \sin(4t) + 4C_2 \cos(4t) \implies y'(0) = 4C_2 = -2 \implies C_2 = -0.5$. \\
    $y(t) = -0.5 \cos(4t) - 0.5 \sin(4t)$. \\
    \textbf{Phase-Amplitude Form:} $A = \sqrt{(-0.5)^2 + (-0.5)^2} = \sqrt{0.25 + 0.25} = \sqrt{0.5} = \frac{\sqrt{2}}{2}$. \\
    $\tan \phi = \frac{C_1}{C_2} = \frac{-0.5}{-0.5} = 1$. Since $C_1 < 0$ ($\sin \phi < 0$) and $C_2 < 0$ ($\cos \phi < 0$), $\phi$ is in Quadrant III. \\
    $\phi = \arctan(1) + \pi = \frac{\pi}{4} + \pi = \frac{5\pi}{4}$. \\
    \textbf{Final Answer:} $y(t) = \frac{\sqrt{2}}{2} \sin\left(4t + \frac{5\pi}{4}\right)$.

    % Solution 2
    \item \textbf{Find Equilibrium Time:} Set $y(t) = 0$. \\
    $\frac{\sqrt{2}}{2} \sin\left(4t + \frac{5\pi}{4}\right) = 0 \implies 4t + \frac{5\pi}{4} = n\pi$. \\
    For the first positive time, choose $n = 2$: \\
    $4t + \frac{5\pi}{4} = 2\pi \implies 4t = \frac{8\pi}{4} - \frac{5\pi}{4} = \frac{3\pi}{4} \implies t = \frac{3\pi}{16}$. \\
    \textbf{Final Answer:} $t = \frac{3\pi}{16}$ seconds.

    % Solution 3
    \item \textbf{Find Period and Frequency:} $m=4, k=25, b=0$. \\
    $\omega = \sqrt{\frac{k}{m}} = \sqrt{\frac{25}{4}} = \frac{5}{2}$ rad/sec. \\
    Period $P = \frac{2\pi}{\omega} = \frac{2\pi}{5/2} = \frac{4\pi}{5}$ seconds. \\
    Frequency $f = \frac{\omega}{2\pi} = \frac{5/2}{2\pi} = \frac{5}{4\pi}$ cycles/sec (Hz). \\
    \textbf{Final Answer:} $P = \frac{4\pi}{5}$ s, $f = \frac{5}{4\pi}$ Hz.

    % Solution 4
    \item \textbf{Set up the IVP:} $m=2, k=18, b=0$. $y(0) = 1, y'(0) = 3$. \\
    $2y'' + 18y = 0 \implies y'' + 9y = 0 \implies r^2 + 9 = 0 \implies r = \pm 3i$. \\
    $y(t) = C_1 \cos(3t) + C_2 \sin(3t)$. \\
    \textbf{Apply Initial Conditions:} $y(0) = C_1 = 1$. \\
    $y'(t) = -3C_1 \sin(3t) + 3C_2 \cos(3t) \implies y'(0) = 3C_2 = 3 \implies C_2 = 1$. \\
    \textbf{Final Answer:} $y(t) = \cos(3t) + \sin(3t)$.

    % Solution 5
    \item \textbf{Phase-Amplitude Form:} $y(t) = \cos(3t) + \sin(3t)$. Here $C_1 = 1, C_2 = 1$. \\
    $A = \sqrt{1^2 + 1^2} = \sqrt{2}$. \\
    $\tan \phi = \frac{C_1}{C_2} = \frac{1}{1} = 1$. Since $C_1 > 0$ and $C_2 > 0$, $\phi$ is in Quadrant I. \\
    $\phi = \arctan(1) = \frac{\pi}{4}$. \\
    \textbf{Final Answer:} $y(t) = \sqrt{2}\sin\left(3t + \frac{\pi}{4}\right)$.

    % Solution 6
    \item \textbf{Set up the IVP:} $m=1, b=5, k=4$. $y(0) = 2, y'(0) = -1$. \\
    $y'' + 5y' + 4y = 0 \implies r^2 + 5r + 4 = 0 \implies (r+1)(r+4) = 0 \implies r = -1, -4$. \\
    $y(t) = C_1 e^{-t} + C_2 e^{-4t}$. \\
    \textbf{Apply Initial Conditions:} \\
    $y(0) = C_1 + C_2 = 2$. \\
    $y'(t) = -C_1 e^{-t} - 4C_2 e^{-4t} \implies y'(0) = -C_1 - 4C_2 = -1$. \\
    Add the two equations: $(C_1 + C_2) + (-C_1 - 4C_2) = 2 - 1 \implies -3C_2 = 1 \implies C_2 = -1/3$. \\
    $C_1 - 1/3 = 2 \implies C_1 = 7/3$. \\
    \textbf{Final Answer:} $y(t) = \frac{7}{3}e^{-t} - \frac{1}{3}e^{-4t}$.

    % Solution 7
    \item \textbf{Find Equilibrium Time:} Set $y(t) = 0$. \\
    $\frac{7}{3}e^{-t} - \frac{1}{3}e^{-4t} = 0 \implies 7e^{-t} = e^{-4t}$. \\
    Divide by $e^{-4t}$: $7e^{3t} = 1 \implies e^{3t} = \frac{1}{7}$. \\
    $3t = \ln(1/7) \implies t = \frac{1}{3}\ln(1/7) \approx -0.648$. \\
    Since $t < 0$, there is no positive time where the mass passes through equilibrium. \\
    \textbf{Final Answer:} The mass never passes through equilibrium for $t > 0$.

    % Solution 8
    \item \textbf{Set up the IVP:} $m=1, b=4, k=4$. $y(0) = 1, y'(0) = 0$. \\
    $y'' + 4y' + 4y = 0 \implies r^2 + 4r + 4 = 0 \implies (r+2)^2 = 0 \implies r = -2$. \\
    $y(t) = (C_1 + C_2 t)e^{-2t}$. \\
    \textbf{Apply Initial Conditions:} $y(0) = C_1 = 1$. \\
    $y'(t) = C_2 e^{-2t} - 2(C_1 + C_2 t)e^{-2t}$. \\
    $y'(0) = C_2 - 2C_1 = 0 \implies C_2 - 2(1) = 0 \implies C_2 = 2$. \\
    \textbf{Final Answer:} $y(t) = (1 + 2t)e^{-2t}$.

    % Solution 9
    \item \textbf{Find Maximum Displacement:} Set $y'(t) = 0$. \\
    From Problem 8, $y'(t) = 2e^{-2t} - 2(1 + 2t)e^{-2t} = e^{-2t}(2 - 2 - 4t) = -4te^{-2t}$. \\
    $-4te^{-2t} = 0 \implies t = 0$. \\
    The maximum displacement occurs at the initial release point. \\
    \textbf{Final Answer:} Maximum displacement is $y(0) = 1$ meter.

    % Solution 10
    \item \textbf{Set up the IVP:} $m=1, b=2, k=5$. $y(0) = 2, y'(0) = -2$. \\
    $y'' + 2y' + 5y = 0 \implies r^2 + 2r + 5 = 0 \implies r = \frac{-2 \pm \sqrt{4 - 20}}{2} = -1 \pm 2i$. \\
    $y(t) = e^{-t}(C_1 \cos(2t) + C_2 \sin(2t))$. \\
    \textbf{Apply Initial Conditions:} $y(0) = C_1 = 2$. \\
    $y'(t) = -e^{-t}(C_1 \cos(2t) + C_2 \sin(2t)) + e^{-t}(-2C_1 \sin(2t) + 2C_2 \cos(2t))$. \\
    $y'(0) = -C_1 + 2C_2 = -2 \implies -2 + 2C_2 = -2 \implies 2C_2 = 0 \implies C_2 = 0$. \\
    \textbf{Final Answer:} $y(t) = 2e^{-t}\cos(2t)$.

    % Solution 11
    \item \textbf{Quasiperiod and Quasifrequency:} From Problem 10, $\beta = 2$. \\
    Quasiperiod $P_q = \frac{2\pi}{\beta} = \frac{2\pi}{2} = \pi$ seconds. \\
    Quasifrequency $f_q = \frac{\beta}{2\pi} = \frac{2}{2\pi} = \frac{1}{\pi}$ cycles/sec. \\
    \textbf{Final Answer:} Quasiperiod is $\pi$ s, quasifrequency is $1/\pi$ Hz.

    % Solution 12
    \item \textbf{Damping Factor:} The damping factor is $Ae^{\alpha t}$. \\
    From Problem 10, $C_1 = 2, C_2 = 0$. $A = \sqrt{2^2 + 0^2} = 2$. $\alpha = -1$. \\
    \textbf{Final Answer:} The damping factor is $2e^{-t}$.

    % Solution 13
    \item \textbf{Find Equilibrium Time:} Set $y(t) = 0$. \\
    $2e^{-t}\cos(2t) = 0$. Since $2e^{-t} \neq 0$, we must have $\cos(2t) = 0$. \\
    $2t = \frac{\pi}{2} + n\pi$. For the first positive time, $n=0$: \\
    $2t = \frac{\pi}{2} \implies t = \frac{\pi}{4}$. \\
    \textbf{Final Answer:} $t = \frac{\pi}{4}$ seconds.

    % Solution 14
    \item \textbf{Determine Damping Type:} $m=3, b=12, k=12$. \\
    Discriminant: $b^2 - 4mk = 12^2 - 4(3)(12) = 144 - 144 = 0$. \\
    \textbf{Final Answer:} Since the discriminant is 0, the system is critically damped.

    % Solution 15
    \item \textbf{Determine Damping Type:} $m=2, b=5, k=3$. \\
    Discriminant: $b^2 - 4mk = 5^2 - 4(2)(3) = 25 - 24 = 1$. \\
    \textbf{Final Answer:} Since $1 > 0$, the system is overdamped.

    % Solution 16
    \item \textbf{Determine Damping Type:} $m=1, b=2, k=10$. \\
    Discriminant: $b^2 - 4mk = 2^2 - 4(1)(10) = 4 - 40 = -36$. \\
    \textbf{Final Answer:} Since $-36 < 0$, the system is underdamped.

    % Solution 17
    \item \textbf{Set up the IVP:} $m=1, b=2, k=2$. $y(0) = 0, y'(0) = 1$. \\
    $r^2 + 2r + 2 = 0 \implies r = -1 \pm i$. \\
    $y(t) = e^{-t}(C_1 \cos t + C_2 \sin t)$. \\
    $y(0) = C_1 = 0 \implies y(t) = C_2 e^{-t}\sin t$. \\
    $y'(t) = C_2(-e^{-t}\sin t + e^{-t}\cos t)$. $y'(0) = C_2(1) = 1 \implies C_2 = 1$. \\
    $y(t) = e^{-t}\sin t$. \\
    \textbf{Find Max Displacement:} Set $y'(t) = 0$. \\
    $-e^{-t}\sin t + e^{-t}\cos t = 0 \implies \sin t = \cos t \implies \tan t = 1 \implies t = \pi/4$. \\
    $y(\pi/4) = e^{-\pi/4}\sin(\pi/4) = \frac{\sqrt{2}}{2}e^{-\pi/4}$. \\
    \textbf{Final Answer:} Maximum displacement is $\frac{\sqrt{2}}{2}e^{-\pi/4}$ meters.

    % Solution 18
    \item \textbf{Set up the IVP:} $m=1, b=6, k=9$. $y(0) = 2, y'(0) = -10$. \\
    $r^2 + 6r + 9 = 0 \implies (r+3)^2 = 0 \implies r = -3$. \\
    $y(t) = (C_1 + C_2 t)e^{-3t}$. \\
    $y(0) = C_1 = 2$. \\
    $y'(t) = C_2 e^{-3t} - 3(C_1 + C_2 t)e^{-3t}$. \\
    $y'(0) = C_2 - 3C_1 = -10 \implies C_2 - 6 = -10 \implies C_2 = -4$. \\
    $y(t) = (2 - 4t)e^{-3t}$. \\
    \textbf{Find Equilibrium Time:} Set $y(t) = 0$. \\
    $2 - 4t = 0 \implies 4t = 2 \implies t = 0.5$. \\
    \textbf{Final Answer:} $t = 0.5$ seconds.

    % Solution 19
    \item \textbf{Find Maximum Displacement:} $y(t) = e^{-t} - e^{-3t}$. Set $y'(t) = 0$. \\
    $y'(t) = -e^{-t} + 3e^{-3t} = 0 \implies e^{-t} = 3e^{-3t}$. \\
    Multiply by $e^{3t}$: $e^{2t} = 3 \implies 2t = \ln 3 \implies t = \frac{1}{2}\ln 3$. \\
    Substitute $t$ back into $y(t)$: \\
    $y\left(\frac{1}{2}\ln 3\right) = e^{-\frac{1}{2}\ln 3} - e^{-\frac{3}{2}\ln 3} = (e^{\ln 3})^{-1/2} - (e^{\ln 3})^{-3/2} = 3^{-1/2} - 3^{-3/2}$. \\
    $= \frac{1}{\sqrt{3}} - \frac{1}{3\sqrt{3}} = \frac{3}{3\sqrt{3}} - \frac{1}{3\sqrt{3}} = \frac{2}{3\sqrt{3}} = \frac{2\sqrt{3}}{9}$. \\
    \textbf{Final Answer:} Maximum displacement is $\frac{2\sqrt{3}}{9}$ meters.

    % Solution 20
    \item \textbf{Theoretical Proof:} Set $y(t) = 0$. \\
    $C_1 e^{r_1 t} + C_2 e^{r_2 t} = 0 \implies C_1 e^{r_1 t} = -C_2 e^{r_2 t}$. \\
    Divide by $C_1 e^{r_2 t}$ (assuming $C_1 \neq 0$, otherwise trivial): \\
    $e^{(r_1 - r_2)t} = -\frac{C_2}{C_1}$. \\
    Since the exponential function $f(t) = e^{kt}$ is strictly monotonic (either always increasing or always decreasing), it can intersect a constant value $-\frac{C_2}{C_1}$ at most exactly one time. Thus, there is at most one real solution for $t$. \\
    \textbf{Final Answer:} The proof is complete by showing the monotonicity of the exponential function.

    % Solution 21
    \item \textbf{Solve Nonhomogeneous ODE:} $y'' + 4y = 3\cos(t)$. \\
    1. Homogeneous: $r^2 + 4 = 0 \implies r = \pm 2i \implies y_h(t) = C_1 \cos(2t) + C_2 \sin(2t)$. \\
    2. Particular: Guess $y_p(t) = A\cos(t) + B\sin(t)$. \\
    $y_p''(t) = -A\cos(t) - B\sin(t)$. \\
    Substitute: $(-A\cos t - B\sin t) + 4(A\cos t + B\sin t) = 3\cos t$. \\
    $3A\cos t + 3B\sin t = 3\cos t \implies 3A = 3 \implies A = 1$, and $3B = 0 \implies B = 0$. \\
    $y_p(t) = \cos(t)$. \\
    \textbf{Final Answer:} $y(t) = C_1 \cos(2t) + C_2 \sin(2t) + \cos(t)$.

    % Solution 22
    \item \textbf{Resonance Form:} $y'' + 16y = 5\sin(4t)$. \\
    Homogeneous roots: $r^2 + 16 = 0 \implies r = \pm 4i$. Natural frequency $\omega = 4$. \\
    The forcing function $5\sin(4t)$ has the same frequency as the natural frequency. \\
    To ensure linear independence from the homogeneous solution, we must multiply the standard guess by $t$. \\
    $y_p(t) = t(A\cos(4t) + B\sin(4t))$. \\
    \textbf{Final Answer:} $y_p(t) = t(A\cos(4t) + B\sin(4t))$. The phenomenon is called \textbf{Pure Resonance}.

    % Solution 23
    \item \textbf{Find the Flaw:} \\
    \textit{Error:} The student used the wrong formula for the discriminant; they added instead of subtracted. \\
    \textit{Correction:} The correct discriminant is $b^2 - 4mk = 16 - 4(2)(4) = 16 - 32 = -16$. Since $-16 < 0$, the system is \textbf{underdamped}.

    % Solution 24
    \item \textbf{Find the Flaw:} \\
    \textit{Error:} The student failed to check the signs of $C_1$ and $C_2$ to place the phase angle $\phi$ in the correct quadrant. \\
    \textit{Correction:} Since $C_1 = -2 < 0$ ($\sin \phi < 0$) and $C_2 = -2 < 0$ ($\cos \phi < 0$), $\phi$ must be in Quadrant III. $\phi = \pi/4 + \pi = 5\pi/4$. \textbf{Final Answer:} $y(t) = 2\sqrt{2}\sin(t + 5\pi/4)$.

    % Solution 25
    \item \textbf{Find the Flaw:} \\
    \textit{Error:} The student forgot to multiply the second term by $t$, which is required for repeated real roots to maintain linear independence. \\
    \textit{Correction:} The correct general solution template for a critically damped system is \textbf{$y(t) = (C_1 + C_2 t)e^{-5t}$}.

    % Solution 26
    \item \textbf{Find the Flaw:} \\
    \textit{Error:} The student set the position $y(t) = 0$ (which finds equilibrium time) instead of setting the velocity $y'(t) = 0$ to find the maximum displacement. \\
    \textit{Correction:} $y'(t) = -e^{-t}\sin(t) + e^{-t}\cos(t) = 0 \implies \tan(t) = 1 \implies t = \pi/4$. \textbf{Final Answer:} $t = \pi/4$ seconds.

    % Solution 27
    \item \textbf{Find the Flaw:} \\
    \textit{Error:} The student simply multiplied the derivatives of the two factors instead of applying the Product Rule $(uv)' = u'v + uv'$. \\
    \textit{Correction:} $y'(t) = (-6e^{-3t})\cos(2t) + (2e^{-3t})(-2\sin(2t))$. \textbf{Final Answer:} $y'(t) = -6e^{-3t}\cos(2t) - 4e^{-3t}\sin(2t)$.

    % Solution 28
    \item \textbf{Set up the IVP:} $m=1, k=2, b=0$. $y(0) = 1, y'(0) = -1$. \\
    $y'' + 2y = 0 \implies r = \pm \sqrt{2}i$. \\
    $y(t) = C_1 \cos(\sqrt{2}t) + C_2 \sin(\sqrt{2}t)$. \\
    $y(0) = C_1 = 1$. \\
    $y'(t) = -\sqrt{2}C_1 \sin(\sqrt{2}t) + \sqrt{2}C_2 \cos(\sqrt{2}t) \implies y'(0) = \sqrt{2}C_2 = -1 \implies C_2 = -1/\sqrt{2}$. \\
    \textbf{Phase-Amplitude Form:} $A = \sqrt{1^2 + (-1/\sqrt{2})^2} = \sqrt{1 + 1/2} = \sqrt{3/2}$. \\
    $\tan \phi = \frac{C_1}{C_2} = \frac{1}{-1/\sqrt{2}} = -\sqrt{2}$. \\
    Since $C_1 > 0$ ($\sin \phi > 0$) and $C_2 < 0$ ($\cos \phi < 0$), $\phi$ is in Quadrant II. \\
    $\phi = \arctan(-\sqrt{2}) + \pi \approx 2.186$ rad. \\
    \textbf{Final Answer:} $y(t) = \sqrt{\frac{3}{2}}\sin(\sqrt{2}t + \arctan(-\sqrt{2}) + \pi)$.

    % Solution 29
    \item \textbf{Set up the IVP:} $m=1, b=2, k=10$. $y(0) = 0, y'(0) = 6$. \\
    $y'' + 2y' + 10y = 0 \implies r^2 + 2r + 10 = 0 \implies r = \frac{-2 \pm \sqrt{4 - 40}}{2} = -1 \pm 3i$. \\
    $y(t) = e^{-t}(C_1 \cos(3t) + C_2 \sin(3t))$. \\
    $y(0) = C_1 = 0 \implies y(t) = C_2 e^{-t}\sin(3t)$. \\
    $y'(t) = C_2(-e^{-t}\sin(3t) + 3e^{-t}\cos(3t))$. \\
    $y'(0) = 3C_2 = 6 \implies C_2 = 2$. \\
    \textbf{Final Answer:} $y(t) = 2e^{-t}\sin(3t)$.

    % Solution 30
    \item \textbf{Theoretical Limit:} Evaluate $\lim_{t \to \infty} (C_1 + C_2 t)e^{-rt}$. \\
    Rewrite as a fraction: $\lim_{t \to \infty} \frac{C_1 + C_2 t}{e^{rt}}$. \\
    As $t \to \infty$, the numerator goes to $\pm \infty$ (if $C_2 \neq 0$) and the denominator goes to $\infty$ (since $r > 0$). This is an $\infty/\infty$ indeterminate form. \\
    Apply L'Hôpital's Rule: take the derivative of the numerator and denominator with respect to $t$. \\
    $\lim_{t \to \infty} \frac{C_2}{r e^{rt}}$. \\
    Since $C_2$ and $r$ are constants and $r > 0$, the denominator $re^{rt} \to \infty$. \\
    Therefore, $\lim_{t \to \infty} \frac{C_2}{r e^{rt}} = 0$. \\
    \textbf{Final Answer:} The limit evaluates to 0, proving the mass eventually settles at equilibrium.

\end{enumerate}

\newpage

% ---------------------------------------------------------
\section{Concept Checklist (Appendix)}
% ---------------------------------------------------------
Use this checklist to review specific concepts and quickly locate the relevant practice problems.

\begin{itemize}
    \item \textbf{Concept 1: Undamped Free Vibrations (IVP \& Equation of Motion)} $\rightarrow$ Questions: 1, 4, 28
    \item \textbf{Concept 2: Amplitude, Period, Frequency, and Phase Angle} $\rightarrow$ Questions: 1, 3, 5, 28
    \item \textbf{Concept 3: Time to First Equilibrium (Undamped)} $\rightarrow$ Questions: 2
    \item \textbf{Concept 4: Damped Free Vibrations - Overdamped} $\rightarrow$ Questions: 6, 15, 19
    \item \textbf{Concept 5: Damped Free Vibrations - Critically Damped} $\rightarrow$ Questions: 8, 14, 18
    \item \textbf{Concept 6: Damped Free Vibrations - Underdamped} $\rightarrow$ Questions: 10, 16, 29
    \item \textbf{Concept 7: Damping Factor, Quasiperiod, Quasifrequency} $\rightarrow$ Questions: 11, 12
    \item \textbf{Concept 8: Maximum Displacement (Optimization)} $\rightarrow$ Questions: 9, 17, 19
    \item \textbf{Concept 9: Time to First Equilibrium (Damped)} $\rightarrow$ Questions: 7, 13, 18
    \item \textbf{Concept 10: Theoretical Proofs \& Limits} $\rightarrow$ Questions: 20, 30
    \item \textbf{Concept 11: Forced Vibrations \& Resonance} $\rightarrow$ Questions: 21, 22
    \item \textbf{Concept 12: Diagnostic Testing ("Find the Flaw")} $\rightarrow$ Questions: 23, 24, 25, 26, 27
\end{itemize}

\end{document}